% Part: computability
% Chapter: computability-theory
% Section: fixed-point-thm

\documentclass[../../../include/open-logic-section]{subfiles}

\begin{document}

\olfileid{cmp}{thy}{fix}
\olsection{Teorema Titik Tetep}

Ayo kita gatekake maneh masalah mandheg. Minangka notasi
sawetara wektu, ayo nulis $\gn{\cfind{x}(y)}$ kanggo
$\tuple{x, y}$; anggepen iki minangka ``jeneng'' kanggo
nilai $\cfind{x}(y)$. Kanthi notasi iki, kita bisa
ngandharake maneh salah siji bukti yen masalah
mandheg ora bisa diputusake.

Pitakon: apa ana fungsi komputabel $h$ kang nduwé
sipat ing ngisor iki? Kanggo saben $x$ lan $y$,
\[
h(\gn{\cfind{x}(y)}) =
\begin{cases}
1 & \text{yen $\cfind{x}(y) \fdefined$} \\
0 & \text{yen ora.}
\end{cases}
\]

Wangsulan: Ora; yen ana, fungsi parsial
\[
g(x) \simeq
\begin{cases}
0 & \text{yen $h(\gn{\cfind{x}(x)}) = 0$} \\
\text{ora ditegesi} & \text{yen ora}
\end{cases}
\]
bakal komputabel, mula nduwé sawijining indeks~$e$.
Nanging banjur kita nduwé
\[
\cfind{e}(e) \simeq
\begin{cases}
0 & \text{yen $h(\gn{\cfind{e}(e)}) = 0$} \\
\text{ora ditegesi} & \text{yen ora,}
\end{cases}
\]
lan ing kahanan iki $\cfind{e}(e)$ ditegesi yen
lan mung yen ora ditegesi, yaiku kontradiksi.

Saiki gatekna persamaan kang ngemot $\cfind{e}$.
Ing kono ana sawijining wujud rujukan marang awake
dhewe: kita wis ngatur supaya nilai $\cfind{e}(e)$
gumantung marang $\gn{\cfind{e}(e)}$ kanthi cara
tartamtu. Teorema titik tetep nyebut yen kita
{\em bisa} nindakake iki sacara umum---ora mung
kanggo mbuktekake kontradiksi.

\olref{lem:fixed-equiv} menehi rong cara ekuivalen
kanggo nyebut teorema titik tetep. Miturut logika,
kalorone ekuivalen amarga kalorone bener; nanging
kang tenan kita maksud yaiku saben pratelan
bisa diturunake kanthi langsung saka liyane,
saengga bisa dadi rong rumusan alternatif
saka teorema kang padha.

\begin{lem}
\ollabel{lem:fixed-equiv}
Pratelan-pratelan ing ngisor iki ekuivalen:
\begin{enumerate}
\item Kanggo saben fungsi komputabel parsial $g(x,y)$,
  ana indeks~$e$ saengga kanggo saben~$y$,
\[
\cfind{e}(y) \simeq g(e,y).
\]
\item Kanggo saben fungsi komputabel~$f(x)$,
  ana indeks~$e$ saengga kanggo saben~$y$,
\[
\cfind{e}(y) \simeq \cfind{f(e)}(y).
\]
\end{enumerate}
\end{lem}

\begin{proof}
$(1) \Rightarrow (2)$: Yen diwenehi $f$, tetepna
$g$ kanthi $g(x,y) \simeq \fn{Un}(f(x),y)$.
Gunakna (1) kanggo ngentukake indeks~$e$
saengga kanggo saben~$y$,
\begin{align*}
\cfind{e}(y) & = \fn{Un}(f(e),y) \\
& = \cfind{f(e)}(y).
\end{align*}

$(2) \Rightarrow (1)$: Yen diwenehi $g$,
gunakna teorema $s$-$m$-$n$ kanggo ngentukake
$f$ saengga kanggo saben $x$ lan~$y$,
$\cfind{f(x)}(y) \simeq g(x,y)$. Gunakna (2)
kanggo ngentukake indeks~$e$ saengga
\begin{align*}
\cfind{e}(y) & = \cfind{f(e)}(y) \\
& = g(e,y).
\end{align*}
Iki ngrampungake bukti.
\end{proof}

\begin{explain}
Sadurunge nuduhake yen pratelan (1) bener
(lan mulane (2) uga bener), pikirna sepira
anehe pratelan iku. Anggepen $e$ minangka
program komputer; pratelan (1) nyebut yen
kanggo fungsi komputabel parsial $g(x,y)$
sembarang, kowe bisa nemokake program
komputer $e$ kang ngitung $g_e(y) \simeq g(e,y)$.
Kanthi tembung liya, kowe bisa nemokake
program komputer kang ngitung fungsi
kang ngrujuk marang program iku dhewe.
\end{explain}

\begin{thm}
Rong pratelan ing \olref{lem:fixed-equiv}
iku bener. Mligine, kanggo saben fungsi
komputabel parsial $g(x,y)$, ana indeks~$e$
saengga kanggo saben~$y$,
\[
\cfind{e}(y) \simeq g(e,y).
\]
\end{thm}

\begin{proof}
Bahan-bahan bukti iki wis kasirat ing rembugan
bab masalah mandheg ing ndhuwur. Tetepna
$\fn{diag}(x)$ minangka fungsi komputabel kang,
kanggo saben $x$, mbalekake indeks kanggo fungsi
$f_x(y) \simeq \cfind{x}(x,y)$, yaiku
\[
\cfind{\fn{diag}(x)}(y) \simeq \cfind{x}(x,y).
\]
Anggepen $\fn{diag}$ minangka fungsi kang ngowahi
program kanggo fungsi rong argumen dadi program
kanggo fungsi siji argumen, kanthi netepake
program asal minangka argumen kapisane.
Fungsi $\fn{diag}$ bisa ditetepake sacara formal
mangkene: luwih dhisik tetepna $s$ kanthi
\[
s(x,y) \simeq \fn{Un}^2(x,x,y),
\]
ing ngendi $\fn{Un}^2$ minangka fungsi telung
argumen kang universal kanggo fungsi komputabel
parsial rong argumen. Banjur miturut teorema
$s$-$m$-$n$, kita bisa nemokake fungsi rekursif
primitif $\fn{diag}$ kang nyukupi
\[
\cfind{\fn{diag}(x)}(y) \simeq s(x,y).
\]

Saiki tetepna fungsi $l$ kanthi
\[
l(x,y) \simeq g(\fn{diag}(x),y).
\]
lan tetepna $\gn{l}$ minangka indeks kanggo $l$.
Pungkasane, tetepna $e = \fn{diag}(\gn{l})$.
Mula kanggo saben $y$, kita nduwé
\begin{align*}
\cfind{e}(y) & \simeq \cfind{\fn{diag}(\gn{l})}(y) \\
& \simeq \cfind{\gn{l}}(\gn{l}, y) \\
& \simeq l(\gn{l}, y) \\
& \simeq g(\fn{diag}(\gn{l}),y) \\
& \simeq g(e, y),
\end{align*}
kaya kang dikarepake.
\end{proof}

\begin{explain}
Apa kang kedadeyan? Upamane kowe diwenehi
tugas nulis program komputer kang nyithak
program iku dhewe. Upamane uga kowe
nganggo basa pamrograman kang nduwé
pustaka fungsi runtutan aksara kang
sugih lan aneh. Mligine, upamane basa
pamrogramanmu nduwé fungsi $\fn{diag}$
kang makarya mangkene: yen diwenehi
runtutan aksara~$s$, $\fn{diag}$ nemokake
saben simbol `x' ing~$s$, banjur ngganteni
simbol iku nganggo versi nganggo tandha
petik saka runtutan aksara asal. Tuladhane,
yen diwenehi runtutan aksara
\begin{quote}
\begin{verbatim}
hello x world
\end{verbatim}
\end{quote}
minangka input, fungsi iku mbalekake
\begin{quote}
\begin{verbatim}
hello 'hello x world' world
\end{verbatim}
\end{quote}
minangka output. Ing kahanan iku,
gampang nulis program kang dikarepake;
kowe bisa mriksa yen
\begin{quote}
\begin{verbatim}
print(diag('print(diag(x))'))
\end{verbatim}
\end{quote}
nindakake iku. Kanggo basa pamrograman
kang luwih umum kaya C++ lan Java,
gagasan kang padha (kanthi cara
ngetrapake kang luwih ruwet) isih bisa.

Saiki kita mung kari sawetara langkah
menyang bukti teorema titik tetep.
Upamane sawijining varian fungsi print
$\fn{print}(x,y)$ nampa runtutan aksara $x$
lan argumen angka liyane $y$, banjur
nyithak runtutan aksara $x$ bola-bali,
cacah $y$ kaping. Mula ``program''
\begin{quote}
\begin{verbatim}
getinput(y);
print(diag('getinput(y); print(diag(x), y)'), y)
\end{verbatim}
\end{quote}
nyithak program iku dhewe $y$ kaping
nalika nampa input $y$. Yen rangka
$\fn{getinput}$---$\fn{print}$---$\fn{diag}$
diganti karo fungsi sembarang $g(x,y)$,
kita ngentukake
\begin{quote}
\begin{verbatim}
g(diag('g(diag(x), y)'), y)
\end{verbatim}
\end{quote}
yaiku program kang, nalika nampa input $y$,
nglakokake $g$ marang program iku dhewe
lan $y$. Yen ``menehi tandha petik''
dianggep minangka ``nganggo indeks kanggo'',
kita ngentukake bukti ing ndhuwur.

Saiki ora apa-apa.\ yen kowe nganggep
bukti iki minangka siasat formal utawa
sihir. Nanging kowe kudune bisa
nyusun maneh rincian argumen ing ndhuwur.
Nalika kita mbuktekake teorema-teorema
ketaklengkapan (lan ``teorema titik tetep''
kang gegayutan), kita bakal ngrembug
cara liyane kanggo mangerteni sebabe
cara iki bisa lumaku.
\end{explain}

\begin{tagblock}{lambda}
\begin{digress}
Gagasan kang padha bisa digunakake kanggo
ngentukake kombinator ``titik tetep''.
Upamane kowe nduwé term lambda $g$,
lan kepengin term liyane $k$ kang
nduwé sipat yen $k$ ekuivalen-$\beta$
karo $gk$. Tetepna term-term
\[
\fn{diag}(x) = xx
\]
lan
\[
l(x) = g(\fn{diag}(x))
\]
miturut konvensi notasi kita; kanthi
tembung liya, $l$ iku term
$\lambd[x][g(xx)]$. Tetepna $k$
minangka term $ll$. Mula kita nduwé
\begin{align*}
k & = (\lambd[x][g(xx)])(\lambd[x][g(xx)]) \\
& \red  g((\lambd[x][g(xx)])(\lambd[x][g(xx)])) \\
& = gk.
\end{align*}
Yen kita njupuk
\[
Y = \lambd[g][((\lambd[x][g(xx)])(\lambd[x][g(xx)]))]
\]
mula $Yg$ lan $g(Yg)$ padha-padha bisa
direduksi menyang sawijining term kang padha;
dadi $Yg \equiv_\beta
g(Yg)$. Iki diarani ``kombinator Curry''.
Yen minangka gantine kita njupuk
\[
Y = (\lambd[xg][g(xxg)])(\lambd[xg][g(xxg)])
\]
mula sejatine $Yg$ bisa direduksi menyang
$g(Yg)$, pratelan kang luwih kuwat.
Versi $Y$ kang kapindho iki diarani
``kombinator Turing''.
\end{digress}
\end{tagblock}

\end{document}
